% !TEX root = ../main.tex
% Introduction
% Claims to support: headline claim + RQs + contributions

\section{Introduction}
\label{sec:01_intro}
Issue Tracking Systems provide a structured way to communicate concerns related to the software project by allowing users to submit issue reports~\cite{colavito2024leveraging}. Issues reports can be classified with labels to help with triage. Effective issue report classification (IRC) results in faster issue resolution~\cite{liao2018exploring}, and is crucial for many software maintenance tasks, such as issue prioritization~\cite{catolino2019not}, issue assignment~\cite{wu2022spatial, jeong2009improving}, and developer onboarding~\cite{santos2023tag, steinmacher2018let, aracena2025applying}. However, studies show that labels are rarely applied in issue creation~\cite{antoniol2008bug, herzig2013s, cabot2015exploring, kallis2019ticket}. As a result, project maintainers must manually classify the issues for triage, which is laborious and time-consuming~\cite{fan2017road}.

To address this problem, numerous automated methods have been proposed around supervised learning over three established issue report categories: \emph{bug}, \emph{feature}, and \emph{question}~\cite{antoniol2008bug, kallis2019ticket, izadi2022catiss, aracena2025applying, aracena2024applyinglargelanguagemodels}. Early work relied on classic machine learning (ML) techniques~\cite{antoniol2008bug, kallis2019ticket, fan2017road, herzig2013s}, while more recent studies fine-tune transformer models on labeled issue reports~\cite{izadi2022catiss, izadi2022predicting, trautsch2022predicting, bharadwaj2022github, colavito2022issue, wang2021well, gomes2023bert}. These approaches require a substantial amount of training data and often struggle with minority classes~\cite{izadi2022catiss, izadi2022predicting, colavito2022issue, heo2025study, vargovich2023givemelabeledissues}. The current state-of-the-art methods fine-tune Large Language Models (LLMs) using Low-Rank Adaptation (LoRA~\cite{hu2022lora})~\cite{heo2025study, aracena2025applying, aracena2024applyinglargelanguagemodels}. However, fine-tuning LLMs is computationally expensive and must be repeated whenever the model is swapped or to incorporate new labeled issues that accumulate in the project~\cite{fan2017road, kallis2019ticket}.

Prior work shows that placing a few demonstration examples in the prompt can boost an LLM's performance on software-related tasks~\cite{assi2026llm, le2023log} through in-context learning~\cite{brown2020language}. Since few-shot prompting enhances performance without model training or weight updates, it has the potential to serve as a cost-efficient alternative to fine-tuning LLMs for IRC. Few-shot performance, however, depends on the choice of in-context examples: retrieval-based example selection outperforms random selection~\cite{liu2022makes}, and the most similar neighbors of a query tend to share the same label, making them ideal classified examples~\cite{cover1967nearest}. The combination of these insights leads to retrieval-augmented few-shot prompting, yet a systematic evaluation of this approach for IRC remains largely underexplored in the literature.

To bridge this gap, we introduce \textbf{R}etrieval-\textbf{A}ugmented \textbf{G}enerative \textbf{TAG}ging (\ragtag). Given a query issue, \ragtag\ retrieves its top-$k$ most textually similar issue reports from an embedded index of existing issue reports and presents them as classified examples in the prompt. $k{=}0$ results in a zero-shot LLM baseline, similar to prior work~\cite{colavito2025benchmarking}. Our evaluation (\Cref{sec:ragtag}) reveals that the \emph{question} class has the weakest \ragtag\ performance, and that question issues are frequently misclassified as bugs. We hypothesize that bug-labeled neighbors retrieved for true-question queries reinforce this misclassification. To tackle this problem, we propose \textbf{B}alanced \ragtag (\bragtag), a \ragtag\ variant that drops the bug-labeled neighbors whenever $N_\text{bug} - N_\text{question} \le m$, where $N_\text{label}$ is the count of neighbors with that label in the top-$k$ and $m$ is a small positive margin.

Following the standard $k$-nearest-neighbor classification framework~\cite{cover1967nearest}, with neighbors weighted by their cosine similarity to the query~\cite{dudani1976distance}, we propose \textbf{Vo}ting-based \textbf{TAG}ging (\votag). Similar to \ragtag\ and \bragtag, for a given query issue, \votag\ first retrieves its top-$k$ neighbors, and predicts the query label by a similarity-weighted vote on their labels. \votag\ establishes a retrieval-only baseline that isolates the contribution of the LLM from that of retrieval in \ragtag\ and \bragtag. It can also be used as a fallback mechanism in case of invalid LLM outputs across all methods due to its low cost.

We evaluate our proposed methods against the established LoRA~\cite{hu2022lora} fine-tuning baseline~\cite{heo2025study, aracena2025applying, aracena2024applyinglargelanguagemodels} across the Qwen2.5-Instruct model family at 3B, 7B, 14B, and 32B parameters~\cite{yang2024qwen25} to analyze model scale effect without architecture variance. We use Heo et al.'s~\cite{heo2025study} 11-project dataset for our experiments and adopt its two evaluation settings: a \emph{project-specific} (PS) configuration where retrieval/training is over each project's individual data split in isolation, and a \emph{project-agnostic} (PA) configuration where retrieval/training is over data from the 11 projects pooled together.

Results show that retrieved few-shot examples consistently improve IRC performance. Even a single retrieved example outperforms the zero-shot baseline across all model sizes (\Cref{sec:ragtag}). LLM reasoning also consistently improves over pure retrieval. Every \ragtag\ configuration outperforms \votag's best result. Notably, however, \votag\ comes within 0.018 macro $F_1$ of Qwen-3B's zero-shot (\Cref{sec:rq1}).

Additionally, retrieval-augmented few-shot prompting is competitive with fine-tuning. Using only each project's own labeled issues (PS), \ragtag\ outperforms fine-tune by 0.021--0.040 macro $F_1$ on all model sizes (\Cref{sec:finetune-eval}). Fine-tune only catches up when training $11\times$ more labeled data via cross-project pooling (PA), and even then \ragtag\ remains competitive: aggregated across all model sizes, fine-tune PA leads \ragtag\ PS by 0.029 macro $F_1$, and the two methods are statistically tied at Qwen-3B and Qwen-32B. \bragtag\ closes this gap by reducing the question-to-bug misclassification rate without sacrificing performance on other labels (\Cref{sec:bragtag}). The aggregate ($\bragtag - $ fine-tune) macro $F_1$ difference is $-0.010$. Applying a \votag\ fallback to invalid LLM outputs across all methods tightens this to statistical equivalence, passing TOST at $\delta{=}0.01$.

Retrieval-augmented few-shot offers deployment advantages as well. Averaged across the four model sizes, \ragtag\ and \bragtag\ require 33\% less peak GPU memory than fine-tuning (\Cref{sec:finetune-eval}). Additionally, incorporating newly labeled issues and swapping the underlying model mid-production do not require retraining cycles.

Overall, this study demonstrates retrieval-augmented few-shot's practicality in IRC, particularly when hardware resources are constrained, labeled data is limited, or the deployment requires frequent model or data updates. 

The remainder of this paper is organized as follows: \Cref{sec:02_related} reviews prior work. \Cref{sec:03_approach} explains our proposed methods and the fine-tuning baseline. \Cref{sec:setup} details the experimental setup. \Cref{sec:evaluations} explains our
evaluations and reports the results. \Cref{sec:discussion} discusses failure modes, practical implications, and future work. \Cref{sec:threats} addresses the potential threats to validity. Finally, \Cref{sec:conclusion} concludes the paper.


